% !PRINTED 687
\weberpage{687}
Die Zahl $T$ ist eine im Allgemeinen nicht genauer zu bestimmende
 Function von $t$, die mit $t$ zugleich ins Unendliche
wächst, die sich aber, da sie nur ganzzahlige Werthe hat, nur
stufenweise, also unstetig ändert. Worauf es hier ankommt, ist
\gesperrt{der Grenzwerth, dem sich das Verhältniss $T:t$ mit
unendlich wachsendem $t$ nähert} (der in dem oben angeführten
 Beispiele $1:m$ ist).

Dieser Grenzwerth lässt sich durch Betrachtungen, wie wir
sie schon im \S.\,155 benutzt haben, finden.

Wir bilden uns zunächst eine Basisform des gegebenen Ideals $\mathfrak{a}$:
\weq{1}{
\alpha = \alpha_1 x_1 + \alpha_2 x_2 + \cdots + \alpha_n x_n,
}
d.\,h. eine lineare Function der $n$ Variablen $x_i$, von der Eigenschaft,
dass durch Substitution aller ganzen rationalen Zahlen für
$x_i$ aus $\alpha$ alle durch $\mathfrak{a}$ theilbare ganze Zahlen des Körpers $\Omega$ entstehen
 (\S.\,146). Mit $\alpha_{i,1}, \alpha_{i,2}, \ldots, \alpha_{i,n}$ bezeichnen wir die mit $\alpha_i$
conjugirten Werthe und bilden die lineare Substitution für $x_i$:
\weq{2}{
\begin{aligned}
y_1 &= \alpha_{1,1} x_1 + \alpha_{2,1} x_2 + \cdots + \alpha_{n,1} x_n\\
y_2 &= \alpha_{1,2} x_1 + \alpha_{2,2} x_2 + \cdots + \alpha_{n,2} x_n\\
&\ \cdot\ \cdot\ \cdot\ \cdot\ \cdot\ \cdot\ \cdot\ \cdot\ \cdot\ \cdot\ \cdot\ \cdot\\
y_n &= \alpha_{1,n} x_1 + \alpha_{2,n} x_2 + \cdots + \alpha_{n,n} x_n.
\end{aligned}
}
Ist $A$ die Determinante dieses Gleichungssystems
\[
A = \Sigma \pm \alpha_{1,1} \alpha_{2,2} \ldots \alpha_{n,n},
\]
so ist $A^2$ die Discriminante des Functionensystems $(\alpha_1, \alpha_2, \ldots, \alpha_n)$,
und daher nach \S.\,147 gleich dem Producte aus dem Quadrate
der absoluten Norm von $\alpha$ und der Körperdiscriminante $\Delta$.
Da die absolute Norm von $\alpha$ zugleich die Norm des Ideals $\mathfrak{a}$
ist (\S.\,151), so setzen wir
\weq{3}{
A^2 = N(\mathfrak{a})^2 \, \Delta.
}
Nun können wir, wie im \S.\,155, die $x_1, x_2, \ldots, x_n$ als
Coordinaten eines Punktes $(x)$ in einem Raume von $n$ Dimensionen
 deuten, und wir lassen jedem Punkt $(x)$ einen Punkt $(x')$
entsprechen, dessen Coordinaten durch
\weq{4}{
x_1' = t^{\frac{1}{n}} x_1, \quad x_2' = t^{\frac{1}{n}} x_2, \quad \ldots, \quad x_n' = t^{\frac{1}{n}} x_n
}
bestimmt sind, worin $t$ eine beliebige positive Grösse bedeutet.
Jedem Gebiete $S$ von Punkten $(x)$ entspricht ein Gebiet $S'$ von
Punkten $(x')$. Alle Figuren in $S$ sind den entsprechenden Figuren
in $S'$ ähnlich, und die Lineardimensionen in $S$ verhalten sich
zu denen in $S'$ wie $1 : t^{\frac{1}{n}}$.
